\begin{lemma}[Outside recoloring and exact exit-leg transport]
Let $G$ be a finite loopless multigraph, with finite vertex and exact edge
types, distinguished edge $e_0$, color set $[k]$, and ordered tree sequence
$T$. Let $\varphi$ be proper and leave exactly $e_0$ uncolored. Let
$\alpha,\theta,\beta$ be pairwise distinct colors, with both $\alpha$
and $\theta$ singleton-closed on $T$ under $\varphi$.
There is an $\alpha,\theta$ \noderef{OutsideKempeRecoloring} $\sigma$
of $T$ from $\varphi$, with selected components avoiding every vertex
of $T$. The coloring $\sigma$ is proper, and exactly $e_0$ is uncolored.
For every ordered tree $U$ with $V(U)\subseteq V(T)$,
$\sigma$ is $(U,C,\varphi)$-stable for every control set $C$, its
tree-missing color set on $U$ equals that of $\varphi$, and every
singleton-closed color on $U$ under $\varphi$ remains singleton-closed
under $\sigma$.
At every outside vertex $a$, missingness of $\alpha$ under $\varphi$
implies missingness of $\theta$ under $\sigma$, and missingness of
$\beta$ is preserved. Moreover, for every ordered tree $A$ with
$V(A)\subseteq V(T)$, every exact $(T,\varphi,\{\alpha,\beta\})$-exit
walk $P$ from $a$ to $b\in V(A)$ becomes an exact
$(A,\sigma,\{\theta,\beta\})$-exit walk $Q$ with the identical ordered
support list and the identical length.
\end{lemma}
\begin{proof}
Put $X=V(T)$. By \noderef{TreeColorSetClosed}, neither an $\alpha$-colored
nor a $\theta$-colored edge crosses the boundary of $X$. Explicitly, if
an affected-color edge $f$ has incident vertices $u\in X$ and $v\notin X$,
then $u\ne v$ and its endpoint pair is $\{u,v\}$, so $f$ belongs to
\noderef{BoundaryEdges} and to the corresponding
\noderef{BoundaryColorEdges}, contrary to singleton closedness.
Consequently the two endpoints of every affected-color edge are on the
same side of $X$.

On the nodes of \noderef{TwoColorIncidenceGraph} for $\alpha,\theta$,
assign a vertex-node the predicate that its vertex is in $X$, and assign
an edge-node the predicate that its exact edge is incident with $X$,
as in \noderef{TreeIncidentEdge}. The preceding endpoint observation
shows that this predicate has equal truth values at adjacent nodes.
Induction on an incidence walk shows that it is constant on every
\noderef{TwoColorComponent}. Select the finite root set
$\{\operatorname{inl}(a):a\notin X\}$. Every vertex in a selected
component is outside $X$, since its root has false predicate.

Define $\sigma=\noderef{OutsideTreeColorSwap}(T,\varphi,\alpha,\theta)$:
fix edges incident with $T$ and exchange the two colors on all other
edges. Every nonincident affected-color edge belongs to a selected
component: choose either endpoint, which is outside $X$, and use its
adjacency to the edge-node. Conversely a selected affected-color edge
cannot be incident with $X$, since an incident endpoint would be in its
selected component. Thus this formula is exactly the exchange on the
selected components, including its treatment of uncolored and unaffected
edges. This verifies every clause of \noderef{OutsideKempeRecoloring}.
Apply \noderef{OutsideKempeRecoloringStable}. It gives properness,
exactly-one-uncoloredness, equality of the colors on every edge incident
with $T$, and equality of all vertex missing sets on $X$.

If $V(U)\subseteq X$, an edge incident with $U$ is incident with $T$.
Therefore the edge clause of \noderef{StableColoring} holds for every
control set, and its vertex clause follows from the preceding equality
of vertex missing sets. By \noderef{TreeMissingColors}, the tree-missing
set is the finite union of these vertex sets; equality of every summand
gives equality of the unions. A boundary edge of $U$ is incident with
$U$, so its color is fixed. The exact boundary-edge set is independent
of coloring; hence a singleton with no colored boundary edge before
the change still has none afterward.

For an outside vertex $a$, suppose $\alpha$ was missing. If a new
$\theta$-edge $f$ were incident with $a$ and were not incident with $X$,
the swap formula would make it an old $\alpha$-edge, a contradiction.
If instead $f$ were incident with $X$, its color would be fixed, so it
would be an old $\theta$-edge with endpoints on opposite sides of $X$,
again a contradiction. This proves the first missingness implication
using \noderef{MissingColors} and \noderef{PresentColors}. No edge gains
or loses color $\beta$, since it differs from both exchanged colors;
this proves the second implication.

Finally let $P$ be the stipulated exit walk. Unfold
\noderef{TreeExitPath}: its support has no vertex of $X$ except $b$,
it is simple, its length is at least two, its starting vertex is outside
$X$, and its starting vertex misses $\alpha$ or $\beta$. Consider an
old $\alpha$-edge node on $P$. It is neither the starting nor the ending
node, since both are vertex-nodes. If the edge were incident with $X$,
the endpoint observation would place both neighboring vertex-nodes on
the walk in $X$. Both would therefore equal $b$, violating simplicity.
Thus every such edge is nonincident and becomes $\theta$. Every old
$\beta$-edge stays $\beta$. Transfer the walk, node by node, into the
new $\theta,\beta$ incidence graph, obtaining $Q$. Its ordered support
and length are unchanged, and it is still simple. Its starting vertex
misses $\theta$ or $\beta$ by the implications just proved and is
outside $A$. Its ending vertex is in $A$ by hypothesis; every other
vertex-node is outside $X$, hence outside $A$. These are all the clauses
of \noderef{TreeExitPath} for $A$, completing the proof.
\end{proof}
